\chapter{Conclusions and Future Directions}
\label{ch:conclusions}
\minitoc

This closing chapter looks back over what
was established and forward to what it opens. The first part gathers the
contributions of the work and weighs them against the objectives set out in
the introduction. The second part maps the directions that the framework,
once understood together with its limits, makes available for further study.

\section{Synthesis of Contributions}
\label{sec:contributions}

The general objective stated in Section~\ref{sec:objectives-scope} was to
provide a self-contained reference on GNNs for homogeneous graphs under the
message-passing framework: a single document that develops the theory, presents
representative architectures, and evaluates them through reproducible
implementations. The work meets this objective, and its specific aims were
addressed in turn across the document.

The foundations that motivate GNNs were established first. Graph-structured
data was characterized as the natural description of relational domains, and
the properties that set it apart from tabular, grid, and sequential data ---
variable size, absence of a canonical ordering, irregular local structure ---
were shown to be precisely the properties that a standard neural network cannot
accommodate. The MLP was examined as the concrete case: its
reliance on a fixed input dimension and a fixed coordinate ordering makes it
unable to operate on graphs without discarding the very structure that carries
the signal. From this failure the requirements on any graph model were derived,
with permutation invariance and equivariance foremost among them.

On these foundations the message-passing framework was developed as the
unifying formulation. A node's representation was defined through the repeated
aggregation of information from its neighborhood, decomposed into the message,
aggregation, and update functions, and shown to satisfy the invariance and
equivariance requirements by construction. The framework was then examined through
its principal implications: the receptive field that successive layers induce, the distinction
between the transductive and inductive settings, the expressive power of the
resulting models and its ceiling, and the design space that the three functions
leave open. The MLP reappeared within this space as the
degenerate case in which the neighborhood is ignored, locating it as a
particular case of the same framework rather than a separate object.

The three representative architectures were then derived as instances of the
framework, one for each level of graph learning. GCN was presented for node
classification, GAE for link prediction through the encoder--decoder view, and
GIN for graph classification, with injective aggregation motivated as the choice that reaches
the expressive ceiling the framework permits. Each architecture was obtained by specializing
the message, aggregation, and update functions, placing the three as points in
the single design space.

The architectures were implemented on standard benchmarks --- Cora for GCN,
Amazon Computers for GAE, and NCI1 for GIN --- in PyTorch Geometric, with
reproducible code in a companion repository. Each was evaluated against a
MLP baseline that assumes no graph structure and predicts from node
features alone. Across all three tasks the architecture that uses the graph
improves on the baseline that does not, confirming in practice the predictive
benefit that the theory anticipates: the relational structure carries
information, and a model built to respect it extracts signal that a
feature-only model cannot.

Finally, the limitations of the framework were examined in their own right.
Oversmoothing, oversquashing, heterophily, and the ceiling on expressive power
were each traced to the framework's defining operation rather than to any
particular architecture, establishing that they are properties of message
passing itself.

The contribution of the work is the reference document itself. Its value lies
not in any single result but in the assembly: a path from the motivation for
learning on graphs through to the limits of the message-passing framework, with the
theory, the architectures, and their evaluation developed in one consistent
notation and one continuous line of argument, intended to serve as a foundational
resource for the study of the field.

\section{Directions for Future Research}
\label{sec:future-work}

The scope of this work was deliberately restricted to homogeneous, undirected,
simple graphs and to three architectures within the message-passing framework.
Its boundaries mark where the field continues, with future research extending
across the limitations of the framework, unexplored regions of its design space, 
new domains and tasks, and architectures that move beyond message passing.

\subsection{Addressing the Limitations of the Framework}
\label{sec:future-limitations}

Chapter~\ref{ch:limitations} established the limitations of message passing as
consequences of its mechanism. Each has since motivated a direction of research
aimed at overcoming it.

The ceiling on useful depth imposed by oversmoothing is addressed by
architectures that give a node access to its earlier, less smoothed
representations rather than only the most recent one, so that depth can grow
without the representations collapsing toward a common
vector~\cite[-5pc]{xu2018}. Oversquashing, the distortion of information from distant 
parts of the graph as it is squeezed through intermediate nodes, is addressed 
by mitigating the associated bottleneck; one direction replaces local propagation 
with global attention mechanisms, a direction discussed separately below. 
Heterophily is met by aggregation schemes that no longer assume neighbors are similar:
separating a node's own representation from its neighbors' and combining
information across several rounds of neighborhood are designs identified as
effective precisely where the homophily assumption fails~\cite[-7pc]{zhu2020}. 
The expressivity ceiling set by the \mbox{1-WL} test is raised by computations that refine
over tuples of nodes rather than single nodes, forming a higher-order hierarchy
that goes beyond this bound~\cite[-3pc]{morris2019}.

\subsection{Unexplored Regions of the Design Space}
\label{sec:future-design-space}

This work examined a limited part of the design space the framework defines.
Other parts of that space contain established architectures it did not develop.

The aggregation step was confined to sum, mean, and max. Replacing it with a
learned, attention-based weighting of neighbors, so that a node determines how much each
neighbor contributes, yields \textit{Graph Attention Networks (GAT)}~\cite[-5pc]{velickovic2018}.
Separating the central node's own representation and learning to aggregate from a
sampled subset of the neighborhood yields \textit{GraphSAGE}, an inductive architecture
built for large graphs and unseen nodes~\cite[-2pc]{hamilton2017}. Both remain within 
message passing while exploring regions of its design space, such as learned aggregation 
and neighborhood sampling.

A further region opens once the graph is no longer static. Where the structure
or the node features evolve in time, the framework can be extended to couple message
passing over the graph with a model of temporal change, as in spatio-temporal architectures 
for forecasting on graphs whose signals evolve over time~\cite{yu2018}.

\subsection{Beyond Homogeneous Graphs and Standard Tasks}
\label{sec:future-domains}

Heterogeneous graphs carry several types of nodes and edges. 
Modeling them requires machinery the homogeneous setting does not: 
transformations specific to each relation type, so that a connection between 
two kinds of entity is treated differently from a connection between 
two others~\cite{schlichtkrull2018}. The homogeneous case developed in this work 
is the prerequisite for these models, which extend message passing to a typed structure.

A different extension changes the unit of prediction. The three tasks of this
work attach predictions to nodes, edges, or whole graphs; a subgraph --- a
portion of a graph with its own internal structure and position within
the larger graph --- is none of these, and learning representations for
subgraphs as primary objects supports tasks 
not captured by the standard levels~\cite{alsentzer2020}. 
Further still, the models considered here are discriminative, predicting properties 
of a given graph; another direction is graph generation, where a model learns a 
distribution over graph structures and samples new instances from it. This 
generative setting has applications in domains such as molecular design~\cite[1pc]{you2018}.

\subsection{Graph Transformers}
\label{sec:future-transformers}

The directions above remain within message passing or extend it. 
A final direction steps outside it. The transformer replaces local neighborhood-based 
propagation in message passing with global attention: every node attends directly 
to every other in a single operation, regardless of distance on the graph~\cite{ying2021}. 
Graph structure, no longer acting as the channel through which information flows, 
is instead incorporated as a positional or structural encoding added to this 
global attention mechanism. This mitigates the bottleneck associated with 
oversquashing --- a distant node is now reachable in a single attention step rather than 
many hops --- at the cost of the locality that gives message passing its scalability and inductive bias.

These directions extend the message-passing framework by relaxing assumptions such as locality,
 expressivity limits, or graph homogeneity, or by transferring the same foundations to 
new domains, tasks, and generative settings. The framework developed in this work is therefore 
not an endpoint but a starting point, providing the basis from which these extensions and 
applications can be understood and evaluated.
